\documentclass{article}
\usepackage[T1]{fontenc}
\usepackage{lmodern}
\usepackage{graphicx}
\usepackage{amsmath,amssymb}
\usepackage[a4paper,margin=2cm]{geometry}
\usepackage[absolute,overlay]{textpos}
\setlength{\TPHorizModule}{1mm}
\setlength{\TPVertModule}{1mm}
\usepackage[indonesian]{babel}
\usepackage{hyperref}
\setlength{\emergencystretch}{2em}
% unit: Halaman 8

\begin{document}

% === Halaman 8 (llm) ===

\begin{itemize}
  \item Jika $a$ adalah akar primitif dari bilangan prima $p$, maka untuk bilangan bulat $b$ kita dapat menemukan pangkat $x$ sedemikian sehingga
  \[ b \equiv a^x \pmod{p}, \quad 0 \le x \le (p-1) \]
  \item Pangkat $x$ disebut \textbf{logaritma diskrit} dari $b$ untuk basis $a \pmod{p}$.
  \item Dalam \textbf{persoalan logaritma diskrit}, diberikan $b \equiv a^x \pmod{p}$, carilah $x$ yang memenuhi kekongruenan tersebut.
  \item Sebagai contoh, 7 adalah akar primitif dari bilangan prima $p = 41$. Maka, carilah $x$ sedemikian sehingga $15 \equiv 7^x \pmod{41}$, jawabannya adalah 3 karena
  \[ 7^3 = 343 \equiv 15 \pmod{41} \]
\end{itemize}

\end{document}
